For a basis blade $e_A = g_{a_1}\cdots g_{a_k}$ ($a_1 < \cdots < a_k$) and a generator $g_b$, $g_be_A = (-1)^{k}e_Ag_b$ if $b \notin A$ and $g_be_A = (-1)^{k-1}e_Ag_b$ if $b \in A$, since $g_b$ anticommutes with every other generator and commutes with itself.

(i) For a bivector blade $g_bg_c$ ($b \ne c$) the two signs multiply: $g_bg_c$ commutes with $e_A$ when $b$, $c$ are both in $A$ or both outside it, and anticommutes when exactly one is in $A$. So conjugation by $g_bg_c$ sends each blade to $\pm$ itself, and an element commutes with $g_bg_c$ if and only if each of its blade components does. A blade $e_A$ with $A$ neither empty nor everything has some $b \in A$ and $c \notin A$ (for $D \ge 2$), and then $g_bg_c$ anticommutes with it. The blades $1$ and $I_D$ commute with every bivector blade, and the bivectors are spanned by the bivector blades. So the commutant is $\mathrm{span}\{1, I_D\}$. The check enumerates every blade for $D = 3, \ldots, 8$ (computed, complete for those $D$).

Clauses (i) to (iii) are standard facts of the classification of real Clifford algebras (the centre and the square of the pseudoscalar; stored as \claimref{primer.sign} and \claimref{std-algebra.pseudoscalar-commutation}; for example Lounesto, \emph{Clifford Algebras and Spinors}, or Porteous); they are proved here in the form used.

(ii) Reversing the product, $I_D\widetilde{I_D} = \prod_a Q(g_a) = (-1)^{D-1}$ and $\widetilde{I_D} = (-1)^{D(D-1)/2}I_D$, so $I_D^2 = (-1)^{D(D-1)/2 + D - 1}$. For $D = 4k, 4k+1, 4k+2, 4k+3$ the exponent's parity is odd, even, even, odd (from $2k(4k-1) + 4k - 1$, $2k(4k+1) + 4k$, $(2k+1)(4k+1) + 4k + 1$, $(4k+3)(2k+1) + 4k + 2$), so $I_D^2 = -1$ exactly for $D \equiv 0, 3 \pmod 4$. The check compares the product with the formula for $D = 2, \ldots, 12$. In the signature $(D-1, 1)$ the product of the squares is $-1$ for every $D$, so $I_D^2 = (-1)^{D(D-1)/2 + 1}$, which is $-1$ exactly for $D \equiv 0, 1 \pmod 4$; the check computes this for $D = 2, \ldots, 12$. The clauses for $D = 5, 6, 7$ therefore depend on the signature: in $\Cl(4, 1)$, $I_5^2 = -1$ and $\pm I_5$ have (a) and (b).

(iii) By (i) an element with (b) is $J = \alpha + \beta I_D$ with real $\alpha$, $\beta$. (c) forces $\alpha = 0$ (a scalar commutes with vectors) and needs $I_D$ to anticommute with every vector, which by the first paragraph ($k = D$, $b \in A$) holds exactly when $D$ is even. Then (a) is $\beta^2I_D^2 = -1$, which needs $I_D^2 = -1$ and $\beta = \pm 1$. By (ii), $D$ even with $I_D^2 = -1$ is $D \equiv 0 \pmod 4$, in either signature (for even $D$ the two formulas agree). For $D = 5, 6$, $I_D^2 = 1$ and $(\alpha + \beta I_D)^2 = \alpha^2 + \beta^2 + 2\alpha\beta I_D \ne -1$ for real $\alpha$, $\beta$, so no element has (a) and (b). For $D = 7$, $(\alpha + \beta I_7)^2 = -1$ forces $\alpha = 0$, $\beta = \pm 1$, and $I_7$, of odd grade in odd dimension, commutes with every generator.

(iv) Moving $g_a$ past the four factors of $I_4$ gives the sign $(-1)^{3}$ for $a \le 3$ (one factor is $g_a$ itself) and $(-1)^{4}$ for $a \ge 4$. A bivector blade $g_bg_c$ picks up the product of its two signs, $-1$ exactly when one index is at most $3$ and the other at least $4$: $4(D-4)$ blades. $I_4^2 = -1$ is (ii) at $D = 4$, unchanged by the extra generators. A rotor $\exp(\theta g_{\mu}g_a/2)$ in a mixed plane then satisfies $I_4\exp(\theta g_{\mu}g_a/2) = \exp(-\theta g_{\mu}g_a/2)I_4$; the check uses the rational rotor $\tfrac35 + \tfrac45 g_1g_4$.

(v) $(L_y \circ R_{I_4})(x) = y(xI_4) = (yx)I_4 = (R_{I_4} \circ L_y)(x)$ by associativity, and $R_{I_4}^2 = R_{I_4^2} = -1$. By (iv), $I_4e_A = (-1)^{k_4}e_AI_4$ with $k_4$ the number of indices of $A$ in $\{0, 1, 2, 3\}$. An even blade has $k_4 + k_{\mathrm{extra}}$ even; those with $k_4$ odd number $\sum_{k_4\ \mathrm{odd}}\binom{4}{k_4}\cdot\sum_{m\ \mathrm{odd}}\binom{D-4}{m} = 8\cdot 2^{D-5} = 2^{D-2}$ for $D > 4$, and none for $D = 4$. The check enumerates every even blade for $D = 4, \ldots, 7$ and confirms the right unit commutes with a mixed rotor acting on the left while the left unit does not.

(vi) By Fact \claimref{field-theory.imported-tower-limit} an adjoined plane gives null $v$, $w$ with $vw + wv = 1$ over the ground; modelled as $v = \tfrac12(g_4 + g_5)$, $w = \tfrac12(g_4 - g_5)$ with $Q(g_4) = 1$, $Q(g_5) = -1$, both orthogonal to the ground. By (iv) $I_4$ commutes with $g_4$ and $g_5$, hence with $v$ and $w$, and anticommutes with $g_{\mu}v$ and $g_{\mu}w$. This is the sense in which the definition's unit commutes with every mode.

The closing sentence follows from (iii) to (v); it concerns continuation through the integer-dimensional algebras $\Cl(1, D-1)$ only. \emph{Argued, not proved:} dimensional regularisation ('t Hooft and Veltman; Breitenlohner and Maison; Fact \claimref{synthesis.imported-dimensional-regularisation-and-the-bmhv-gamma5}, whose split of $\gamma_5$ is $\{\gamma_5, \bar{\gamma}^{\mu}\} = 0$, $[\gamma_5, \hat{\gamma}^{\mu}] = 0$) is a formal continuation to non-integer $D$, where the same split is the known $\gamma_5$ obstruction; by (ii) the properties have period four in $D$, so none of the algebraic continuations interpolates all three. For vector-like couplings the right unit of (v) behaves as the standard scheme's external $i$, so one-loop quantum electrodynamics can be dimensionally regularised with no input beyond the standard scheme; the one-loop result \claimref{synthesis.one-loop-qed-in-the-spacetime-algebra}, regularised by Pauli-Villars, likewise needs the unit outside the Clifford product on massive lines, which is this right unit. For chiral couplings the series uses (c) to read one chirality as complex-linear and the other as conjugated (Proposition \claimref{field-theory.second-quantisation}); on the mixed blades of (v) that identification fails exactly as $\gamma_5$'s anticommutation fails in the scheme, so chiral loops inherit its finite symmetry-restoring counterterms. Regulators that stay in four dimensions (Pauli-Villars, the lattice, the Epstein-Glaser extension of distributions, and the tower's floors, a sharp mode cut-off that is not Lorentz covariant, \claimref{synthesis.floor-regulator-keeps-the-charge-and-the-central-term-below-the-cut-off}) need no continuation of the unit. Transport-absent.
